% Unit 048: English translation of chapter4.tex, lines 9--60.
% Source: Methods of Algebra, Volume 2, commit
% 9a5803ff77dd3257484cb177f851a73770a59dd3 (CC BY 4.0).
% stable-unit-id: o014.aljabr2.chapter4.overview
% Coverage: Chapter 4 overview before Section 4.1.

% segment-id: o014.li2.chapter4.overview.h001
\chapter{Triangulated Categories and Derived Categories\texorpdfstring{\hspace{0.1em}}{}}%
\label{sec:triangulated-derived-cat}
\hypertarget{o014.aljabr2.chapter4.overview}{ }

% segment-id: o014.li2.chapter4.overview.p001
Derived categories were pioneered by A.\ Grothendieck and J.-L.\ Verdier in
the 1960s to study duality theorems in algebraic geometry. Subsequent
developments have shown that they provide a convenient and powerful language
in geometry, topology, and representation theory alike. The information
recorded by the derived category $\cate{D}(\mathcal{A})$ of an Abelian
category $\mathcal{A}$ lies between a complex and its cohomology. It is
expressed precisely in the language of triangulated categories introduced in
\S\S~\sourcecrossref{sec:triangular-def}{4.1}--%
\sourcecrossref{sec:triangular-basic}{4.2}: $\cate{D}(\mathcal{A})$ is
equipped with the translation autoequivalence $X\mapsto X[1]$ and a class of
sequences of morphisms called distinguished triangles (also called exact
triangles),
% segment-id: o014.li2.chapter4.overview.q001
\[
  X \xrightarrow{f} Y \xrightarrow{g} Z \xrightarrow{h} X[1],
  \quad \text{or, for short,}\quad X\to Y\to Z\xrightarrow{+1} .
\]
These triangles satisfy axioms (TR0)--(TR5). For example, a distinguished
triangle can be rotated into
% segment-id: o014.li2.chapter4.overview.q002
\[
  Y \xrightarrow{g} Z \xrightarrow{h} X[1] \xrightarrow{-f[1]} Y[1]
  \quad\text{or}\quad
  Z[-1] \xrightarrow{-h[-1]} X \xrightarrow{f} Y \xrightarrow{g} Z .
\]

% segment-id: o014.li2.chapter4.overview.p002
Distinguished triangles replace short exact sequences in the category of
complexes $\cate{C}(\mathcal{A})$, though the more accurate analogy is with
cofiber sequences in homotopy theory. Cohomological functors from a
triangulated category to an Abelian category provide a mechanism for
producing long exact sequences from distinguished triangles.

% segment-id: o014.li2.chapter4.overview.p003
The passage from $\cate{C}(\mathcal{A})$ to $\cate{D}(\mathcal{A})$ is the
subject of \S~\sourcecrossref{sec:derived-cat}{4.4} and proceeds in two
steps.
% segment-id: o014.li2.chapter4.overview.l001
\begin{enumerate}
	% segment-id: o014.li2.chapter4.overview.i001
	\item The first step passes from $\cate{C}(\mathcal{A})$ to
	$\cate{K}(\mathcal{A})$. The latter is equipped with the translation
	functor $X\mapsto X[1]$ and the structure of a triangulated category. For
	every morphism $f:X\to Y$, its mapping cone together with the canonical
	morphisms
	% segment-id: o014.li2.chapter4.overview.q003
	\[
	  X \xrightarrow{f} Y \xrightarrow{\alpha(f)} \Cone(f)
	  \xrightarrow{\beta(f)} X[1]
	\]
	forms the corresponding distinguished triangle, up to isomorphism. The
	tools needed for this step were prepared in the first half of
	\CHref{sec:cplx}.
	% segment-id: o014.li2.chapter4.overview.i002
	\item The second step formally adjoins inverses to all
	quasi-isomorphisms, or equivalently, formally makes every acyclic complex
	zero, by localization. The result is $\cate{D}(\mathcal{A})$, which has
	the same collection of objects as $\cate{C}(\mathcal{A})$ and inherits
	its triangulated structure from $\cate{K}(\mathcal{A})$. This step rests
	on the general theory of Verdier localization in
	\S~\sourcecrossref{sec:triangulated-localization}{4.3}. It is here that
	the octahedral axiom (TR5), the most intricate axiom in the definition of
	a triangulated category, plays its role.
\end{enumerate}

% segment-id: o014.li2.chapter4.overview.p004
Taking cohomology of complexes gives the cohomological functor
$\Hm^0:\cate{D}(\mathcal{A})\to\mathcal{A}$ at the level of the derived
category, with $\Hm^n(X):=\Hm^0(X[n])$. Starting instead from the category
of bounded-below (or bounded-above, or bounded) complexes, the same procedure
gives the triangulated category $\cate{D}^+(\mathcal{A})$ (or
$\cate{D}^-(\mathcal{A})$, or $\cate{D}^{\bdd}(\mathcal{A})$). These may
also be regarded as full triangulated subcategories of
$\cate{D}(\mathcal{A})$, characterized respectively by
$n\ll0$ (or $n\gg0$, or $|n|\gg0$) $\implies\Hm^n(X)=0$.

% segment-id: o014.li2.chapter4.overview.p005
More generally, cohomology-vanishing conditions can be used to define the
full subcategories $\cate{D}^{\geq0}(\mathcal{A})$ and
$\cate{D}^{\leq0}(\mathcal{A})$ of $\cate{D}(\mathcal{A})$, and the
truncation functors on complexes can still be defined on the derived
category. The canonical functor
$\mathcal{A}\to\cate{D}(\mathcal{A})$ is fully faithful and identifies
$\mathcal{A}$ with
$\cate{D}^{\leq0}(\mathcal{A})\cap\cate{D}^{\geq0}(\mathcal{A})$. If
$\mathcal{A}$ has enough injective or projective objects, then
$\Hom_{\cate{D}(\mathcal{A})}(X,Y[n])\simeq
\Ext^n_{\mathcal{A}}(X,Y)$, where $X$ and $Y$ are objects of $\mathcal{A}$
viewed as complexes concentrated in degree zero.%
\footnote{Translator's note (O014-C071): the source describes $X$ and $Y$ as
complexes over $\mathcal{A}$. The fully faithful embedding
$\mathcal{A}\to\cate{D}(\mathcal{A})$ in the preceding sentence and the
definition of $\Ext^n_{\mathcal{A}}$ in \S~\ref{sec:Ext-Tor}, however, give
this formula for objects of $\mathcal{A}$ viewed as degree-zero complexes;
the intended scope has been restored here.}
The definition of $\Ext^n_{\mathcal{A}}$ is given in
\S~\ref{sec:Ext-Tor}, and all these matters are treated in
\S~\sourcecrossref{sec:Hom-Ext}{4.5}.

% segment-id: o014.li2.chapter4.overview.p006
For a left exact functor $F:\mathcal{A}\to\mathcal{A}'$, its classical right
derived functor is realized at the level of bounded-below derived categories
as $\mathrm{R}F:\cate{D}^+(\mathcal{A})\to\cate{D}^+(\mathcal{A}')$. It
can be characterized as a left Kan extension; the relevant data are expressed
by the following $2$-cell diagram.
% segment-id: o014.li2.chapter4.overview.g001
\[\begin{tikzcd}
	\cate{K}^+(\mathcal{A}) \arrow[r, "{\cate{K}F}"] \arrow[d, "Q"'] &
	\cate{K}^+(\mathcal{A}') \arrow[d, "{Q'}"] \arrow[Rightarrow, ld] \\
	\cate{D}^+(\mathcal{A}) \arrow[r, "{\mathrm{R}F}"'] &
	\cate{D}^+(\mathcal{A}')
\end{tikzcd}\]

Here $Q$ and $Q'$ denote the localization functors, while
$\mathrm{R}F$ must also be equipped with a triangulated-functor structure.
The classical $\mathrm{R}^nF$ is obtained by restricting $\mathrm{R}F$ to
$\mathcal{A}$ and then taking $\Hm^n$. If instead $F$ is right exact, the
lift of the classical left derived functor $\mathrm{L}_nF$ is represented as
follows.%
\footnote{Translator's note (O014-C070): in the following diagram, the source
labels the right vertical arrow $Q$, even though its source and target are
$\cate{K}^-(\mathcal{A}')$ and $\cate{D}^-(\mathcal{A}')$, respectively.
In accord with the preceding diagram and the names used in the prose, the
label $Q'$ has been restored here.}
% segment-id: o014.li2.chapter4.overview.g002
\[\begin{tikzcd}
	\cate{K}^-(\mathcal{A}) \arrow[r, "{\cate{K}F}"] \arrow[d, "Q"'] &
	\cate{K}^-(\mathcal{A}') \arrow[d, "{Q'}"] \\
	\cate{D}^-(\mathcal{A}) \arrow[r, "{\mathrm{L}F}"']
	\arrow[ru, Rightarrow] & \cate{D}^-(\mathcal{A}').
\end{tikzcd}\]

% segment-id: o014.li2.chapter4.overview.p007
As in the classical theory, the existence of derived functors and the
detailed structure of derived categories depend on resolutions of complexes.
Resolutions provide the scaffolding needed to work concretely with derived
categories. Details appear in
\S\S~\sourcecrossref{sec:triangulated-functor-localization}{4.6}--%
\sourcecrossref{sec:bounded-derived-functor}{4.8}. The unbounded version
$\cate{D}(\mathcal{A})\to\cate{D}(\mathcal{A}')$ has a similar formulation
and will be discussed in
\S~\sourcecrossref{sec:unbounded-derived}{4.11}.

% segment-id: o014.li2.chapter4.overview.p008
By considering the amplitude of a derived functor, that is, the “shift” it
produces at the level of cohomology, one can discuss the dimension of a
derived functor
(Definition--Proposition~\sourcecrossref{def:dim-functor}{in \S~4.7}).
Under suitable conditions, taking derived functors is also compatible with
composition
(Theorem~\sourcecrossref{prop:localization-injective}{in \S~4.8}). This
provides a uniform explanation for several spectral sequences in the
classical theory, called Grothendieck spectral sequences and described in
Theorem~\sourcecrossref{prop:Grothendieck-ss}{in \S~4.8}. The general theory
of spectral sequences is deferred to
Chapter~\sourcecrossref{sec:ss}{5}.

% segment-id: o014.li2.chapter4.overview.p009
All these constructions can be formulated for general triangulated categories
and their Verdier localizations. It is worth noting that applications often
involve triangulated categories that do not arise directly from Abelian
categories, or cohomological functors other than $\Hm^n$. In this general
setting, the roles of truncation functors and of
$\cate{D}^{\leq0}$ and $\cate{D}^{\geq0}$ are taken by a $t$-structure on
the triangulated category, but this book does not treat that topic.

% segment-id: o014.li2.chapter4.overview.p010
The derived functors discussed concretely in this chapter are limited to
examples arising in algebra. The first is the functor $\RHom$ in
\S~\sourcecrossref{sec:RHom}{4.9}. This is the derived-category version of
the $\Ext$ functor and involves the theory of derived bifunctors. Several
adjunctions will also be lifted to derived categories, following the principle
of replacing $\Hom$ by $\RHom$. Next,
\S~\sourcecrossref{sec:Rlim}{4.10} treats products and coproducts in
$\cate{D}(\mathcal{A})$ and then introduces the lift of $\lim^n$, denoted
$\mathrm{R}\lim$. Like $\lim^n$, this functor has a concise description in
terms of the homotopy limit
(Definition~\sourcecrossref{def:holim}{in \S~4.10}), though at the level of
derived categories the description can only be given in a coarse form.

% segment-id: o014.li2.chapter4.overview.p011
Then \S~\sourcecrossref{sec:otimesL}{4.12} introduces the derived-category
version of the $\Tor$ functor, namely the derived tensor product $\otimesL$,
a triangulated bifunctor. That section also proves many facts, including the
associativity constraint
(Proposition~\sourcecrossref{prop:otimesL-associativity-constraint}{in
\S~4.12}) and the $\RHom$ version of the adjunction
(Theorem~\sourcecrossref{prop:RHom-otimesL-adjunction}{in \S~4.12}). As in
the classical theory, $\RHom$ and $\otimesL$ are the two fundamental examples
of derived functors. A complete and clean description requires unbounded
derived categories and the derived functors between them. This in turn
requires the theory of
\S~\sourcecrossref{sec:unbounded-derived}{4.11}, especially the notions of
K-projective and K-flat resolutions. In the module-theoretic setting
considered here, the existence of these resolutions poses no problem.

% segment-id: o014.li2.chapter4.overview.p012
Finally, several shortcomings of the derived category
$\cate{D}(\mathcal{A})$ should be noted.
% segment-id: o014.li2.chapter4.overview.l002
\begin{enumerate}
	% segment-id: o014.li2.chapter4.overview.i003
	\item Among the axioms of a triangulated category, (TR2) says that every
	morphism can be placed in a distinguished triangle. Such a triangle is
	unique up to isomorphism, but not “up to unique isomorphism” in the sense
	familiar from algebra.
	% segment-id: o014.li2.chapter4.overview.i004
	\item Triangulated categories usually have neither kernels nor cokernels,
	making the operations $\varinjlim$ and $\varprojlim$ extremely difficult
	to carry out. Mapping cones provide substitutes in the homotopical sense,
	but, as already noted, their uniqueness in the derived category is
	problematic.
	% segment-id: o014.li2.chapter4.overview.i005
	\item Although derived functors at the level of derived categories have a
	clear definition, unlike universal cohomological $\delta$-functors they
	are difficult to characterize simply by effaceability; see
	Proposition~\ref{prop:erasable-univ-delta}.
\end{enumerate}

% segment-id: o014.li2.chapter4.overview.p013
The root of the problem is that the homotopy relation is erased too crudely
when $\cate{D}(\mathcal{A})$ is constructed. These shortcomings must be
remedied by a construction that is higher than the derived category and
perhaps also more natural; the derived category should be only the shadow
cast by that construction.

% segment-id: o014.li2.chapter4.overview.n001
\begin{wenxintishi}
	Most of this chapter consists of core concepts in the theory of
	triangulated and derived categories. Side topics include the treatment of
	$n$-extensions in \S~\sourcecrossref{sec:Hom-Ext}{4.5}, centered mainly
	on the theorem of Nobuo Yoneda
	(Theorem~\sourcecrossref{prop:Yoneda-Ext}{in \S~4.5}), and the treatment
	of unbounded derived functors in
	\S~\sourcecrossref{sec:unbounded-derived}{4.11}, which rests on the theory
	of \S~\ref{sec:K-injectives}. Both topics are common knowledge among
	specialists, but readers should judge for themselves whether they need
	them. If one considers only derived functors between bounded-below or
	bounded-above derived categories, every statement in
	\S~\sourcecrossref{sec:otimesL}{4.12} concerning K-projective (or K-flat)
	resolutions may be replaced by a projective (or flat) resolution by
	bounded-above complexes without disrupting the exposition. The validity
	of this replacement is guaranteed by
	Proposition~\sourcecrossref{prop:flat-K-flat}{in \S~4.11}.
\end{wenxintishi}
